\section{Expected Challenges and Mitigations}

Four principal challenges are anticipated, each with a concrete mitigation strategy.

\begin{itemize}

\item \textbf{Scale mismatch.} GRACE FO has a native resolution of approximately 300\nobreakspace km versus Denmark's 43,000\nobreakspace km\textsuperscript{2}. Mitigations: sub-monthly L1B assimilation (Retegui-Schiettekatte \textit{et al.}\nobreakspace 2025), \textbf{ConBay} merging with SMAP (Mehrnegar \textit{et al.}\nobreakspace 2023), ML-based downscaling against Jupiter wells, and \textbf{GIA} correction via ICE-6G\_D.

\item \textbf{Equifinality.} Multiple parameter sets can reproduce equivalent streamflow but incorrect partitioning between evapotranspiration, soil moisture, and recharge. Mitigations: GEUS-derived priors on aquifer parameters, localisation and inflation within \textbf{LETKF}, and \textbf{dPL} as deterministic precursor before stochastic assimilation.

\item \textbf{ML physical consistency.} Pure LSTMs may extrapolate poorly and do not guarantee mass conservation. Mitigations: keeping VIC as the physical backbone, restricting the LSTM to a residual post-processor role, and applying conformal prediction for honest uncertainty intervals.

\item \textbf{Index collapse.} Naive averaging of indices with different time scales loses causal propagation and tail dependence of extremes. Mitigations: pillar-appropriate scales (SPI 1--3 months, SGI 6--12 months), copula aggregation preserving tail dependence (Hao and AghaKouchak 2013), and validation against Naturskaderådet claims, crop yield anomalies, and EDO CDI v4.

\end{itemize}
